%%%%%%%%%%%%%%%%%%%%%%%%%%%%%%%%%%%%%%%%%%%%%%%%%%%%%%%%%%%%
% GLOBAL MACROS BY ENOCH YU
%%%%%%%%%%%%%%%%%%%%%%%%%%%%%%%%%%%%%%%%%%%%%%%%%%%%%%%%%%%%

% Boxes
\usepackage{tcolorbox}
\tcbuselibrary{theorems}
\newcounter{boxcounter}[section]
\renewcommand{\theboxcounter}{\thesection.\arabic{boxcounter}}

\newtcbtheorem[
    auto counter,
    number within=section,
    use counter=boxcounter
]{theostyle}{Theorem}{%
    colback=black!5,
    colframe=black,
    fonttitle=\bfseries,
    sharp corners
}{theo}

\newtcbtheorem[
    auto counter,
    number within=section,
    use counter=boxcounter
]{lemstyle}{Lemma}{%
    colback=Sepia!5,
    colframe=Sepia,
    fonttitle=\bfseries,
    sharp corners
}{lem}

\newtcbtheorem[
    auto counter,
    number within=section,
    use counter=boxcounter
]{probstyle}{Exercise}{%
    colback=BrickRed!3,
    colframe=red!40!black,
    fonttitle=\bfseries,
    sharp corners
}{probstyle}

\newtcbtheorem[
    auto counter,
    number within=section,
    use counter=boxcounter
]{corstyle}{Corollary}{%
    colback=blue!3,
    colframe=blue!30!black,
    fonttitle=\bfseries,
    sharp corners
}{cor}

\newtcbtheorem[
    auto counter,
    number within=section,
    use counter=boxcounter
]{defstyle}{Definition}{%
    colback=yellow!10,
    colframe=Sepia,
    fonttitle=\bfseries,
    sharp corners
}{def}

\newenvironment{theorem}
{\begin{theostyle}}{\end{theostyle}}
\newenvironment{corollary}
{\begin{corstyle}}{\end{corstyle}}
\newenvironment{lemma}
{\begin{lemstyle}}{\end{lemstyle}}
\newenvironment{definition}
{\begin{defstyle}}{\end{defstyle}}
\newenvironment{problem}
{\begin{probstyle}}{\end{probstyle}}

\usepackage{mdframed}
\renewmdenv[
    topline=false,
    bottomline=false,
    rightline=false,
    linecolor=black,
    linewidth=2pt,
    skipabove=\medskipamount,
    skipbelow=\medskipamount,
    frametitle=\it Proof.
]{proof}

\newmdenv[
    topline=false,
    bottomline=false,
    rightline=false,
    linecolor=BrickRed,
    linewidth=2pt,
    skipabove=\medskipamount,
    skipbelow=\medskipamount,
    frametitle=Solution.
]{solution}

\newmdenv[
    backgroundcolor=Periwinkle!3,
    topline=false,
    bottomline=false, 
    rightline=false,
    linecolor=RoyalPurple,
    linewidth=2pt,
    skipabove=\medskipamount, 
    skipbelow=\medskipamount,
    frametitle=In my head...
]{thoughts}

\newmdenv[
    backgroundcolor=yellow!3,
    topline=false,
    bottomline=false,
    rightline=false,
    linecolor=BrickRed,
    linewidth=2pt,
    skipabove=\medskipamount,
    skipbelow=\medskipamount,
    fontcolor=BrickRed,
    frametitle=Tip.
]{tip}



% Example
\NewDocumentEnvironment{example}{o}
{
 \IfValueT{#1}{\vspace{11pt}}
 \begin{mdframed}[
    linecolor=CadetBlue,
    linewidth=1pt,
    skipabove=\medskipamount,
    skipbelow=\medskipamount,
    backgroundcolor=cyan!5,
    innertopmargin=10pt,
    innerbottommargin=10pt,
    innerleftmargin=10pt,
    innerrightmargin=10pt
 ]
 \IfValueT{#1}{%
    \vspace{-11pt}%
    \noindent
    \begin{tikzpicture}[remember picture, overlay]
        \node[
            draw=CadetBlue,
            line width=1pt,
            fill=cyan!5,
            text=Sepia,
            font=\bfseries\small,
            anchor=west,
            xshift=2pt
        ] {\,\strut #1\,};
    \end{tikzpicture}\par
    \vspace{5pt}%
 }%
}
{\end{mdframed}}


% Code
\usepackage{listings}
\definecolor{bg}{RGB}{231,234,234}
\definecolor{kw}{RGB}{0,64,153}
\definecolor{str}{RGB}{163,21,21}
\definecolor{cmt}{RGB}{100,100,100}

\lstdefinestyle{basic}{
    backgroundcolor=\color{bg},
    basicstyle=\ttfamily\small,
    numbers=none,
    frame=none,
    showstringspaces=false,
    tabsize=4,
    showtabs=false,
    breaklines=true,
    aboveskip=0pt,
    belowskip=0pt
}

\lstdefinestyle{cppstyle}{
    language=C++,
    backgroundcolor=\color{bg},
    basicstyle=\ttfamily\small,
    keywordstyle=\color{kw}\bfseries,
    stringstyle=\color{str},
    commentstyle=\color{cmt}\itshape,
    numberstyle=\tiny\color{gray},
    numbers=left, numbersep=5pt,
    breaklines=true,
    rulecolor=\color{gray!50},
    showstringspaces=false,
    tabsize=4, captionpos=b,
    linewidth=\textwidth,
    aboveskip=0pt,
    belowskip=0pt
}

\lstdefinestyle{jin}{
    language=Python,
    backgroundcolor=\color{bg},
    basicstyle=\ttfamily\small,
    keywordstyle=\color{kw}\bfseries,
    stringstyle=\color{str},
    commentstyle=\color{cmt}\itshape,
    numberstyle=\tiny\color{gray},
    numbers=left, numbersep=8pt,
    breaklines=true,
    frame=single,
    rulecolor=\color{gray!50},
    showstringspaces=false,
    tabsize=4, captionpos=b,
    xleftmargin=1.2em,
    framexleftmargin=1.2em,
    belowskip=0pt,
    morekeywords={
        model,compile,fit,predict,activation,
        units,optimizer,loss,epochs,lr,np,plt
    }
}

\lstdefinestyle{jout}{
    basicstyle=\ttfamily\small,
    breaklines=true,
    showstringspaces=false,
    tabsize=4, captionpos=b,
    aboveskip=4pt
}

\newtcbtheorem[no counter]{basicboxstyle}{}{%
    titlebox=invisible,
    detach title,
    colframe=bg,
    colback=bg,
    top=0pt, bottom=0pt, right=0pt, left=0pt,
    boxrule=0pt,
    sharp corners
}{cpp}
\newenvironment{basicbox}
{\begin{basicboxstyle}{}{}}{\end{basicboxstyle}}

\newtcbtheorem[
    auto counter,
    number within=section,
    use counter=boxcounter
]{cppboxstyle}{Code}{%
    colframe=Sepia!80,
    colback=bg,
    top=0pt, bottom=0pt, right=0pt, left=0pt,
    boxrule=0pt,
    sharp corners
}{cpp}
\newenvironment{cppbox}
{\begin{cppboxstyle}}{\end{cppboxstyle}}

\newtcbtheorem[
    auto counter,
    number within=section,
    use counter=boxcounter
]{outputboxstyle}{Output}{%
    colframe=black,
    colback=bg,
    top=0pt, bottom=0pt, right=0pt, left=0pt,
    boxrule=0pt,
    sharp corners
}{cpp}
\newenvironment{outputbox}
{\begin{outputboxstyle}{}{}}{\end{outputboxstyle}}

\newenvironment{jupyter}{
\begin{mdframed}[
    backgroundcolor=white,
    linecolor=Sepia!80,
    linewidth=1pt,
    innerleftmargin=10pt,
    innerrightmargin=10pt,
    innertopmargin=20pt,
    innerbottommargin=0pt,
    skipabove=10pt,
    skipbelow=10pt
]}{\end{mdframed}}



\DeclareMathOperator{\Span}{span}
\DeclareMathOperator{\Row}{row}
\DeclareMathOperator{\Col}{col}
\DeclareMathOperator{\Rank}{r}
\DeclareMathOperator{\Null}{null}
\DeclareMathOperator{\Tr}{tr}



%%%%%%%%%%%%%%% END %%%%%%%%%%%%%%%



